\documentclass[a4paper,10pt]{article}

\usepackage[utf8]{inputenc}
\usepackage[T1]{fontenc}
\usepackage{geometry}
\usepackage{graphicx}
\usepackage{amsmath, amssymb}
\usepackage{booktabs, array, tabularx, longtable}
\usepackage[table]{xcolor}
\usepackage{hyperref}
\usepackage{enumitem}
\usepackage{csquotes}
\usepackage{listings}
\usepackage{footmisc}
\usepackage{graphicx}

\usepackage{algorithm}
\usepackage{algpseudocode}


% Margins
\geometry{margin=2.5cm}

% Color definitions for DONE/TODO
\definecolor{donegreen}{RGB}{34,139,34}
\definecolor{todored}{RGB}{178,34,34}
\newcommand{\done}{\textcolor{donegreen}{\bfseries DONE}}
\newcommand{\todo}{\textcolor{todored}{\bfseries TODO}}

% Fixed centered column for tabularx
\newcolumntype{C}{>{\centering\arraybackslash}X}

% Bibliography
\usepackage[backend=biber,style=alphabetic,sorting=nyt]{biblatex}
\addbibresource{DocSmulatorBib.bib}

\title{\textbf{RMFS Simulator Documentation}}
\author{}
\date{}

\begin{document}
\maketitle
\tableofcontents
\newpage

% ─────────────────────────────────────────────────────────────────────────────
\section{Assumptions}

For the sake of simplicity, the simulator is based on the following simplifying assumptions.

\begin{itemize}[noitemsep]
    \item Replenishment missions are not considered, therefore no inventory constraint is modelled. Each pod contains a limited set of SKUs, but the exact quantity or volume of each item is not known.
    \item Technical aspects of the robots, such as energy management and low-level routing, are outside the scope of this simulator and are not modelled.
    \item At the end of a mission, the pod will return to its storage position.
\end{itemize}

% ─────────────────────────────────────────────────────────────────────────────
\section{General Description}
The discrete-event simulator mimics the operational dynamics of an Amazon-style RMFS (Robotic Mobile Fulfillment System), based on a \textit{parts-to-picker} approach: robots bring pods to workstations, where operators pick items and process orders.

The simulator supports two decision-making modes:
\begin{enumerate}[noitemsep]
    \item \textbf{Cyclic optimisation:} every $\Delta t_{\text{opt}}$ a mathematical model is solved for task generation and sequencing.
    \item \textbf{Greedy policies:} online decisions are made based on policies with limited lookahead.
\end{enumerate}

In both modes, the assignment of a task to a robot is always managed via a greedy policy based on the minimization of the distance between the resting position of the robot and the storage location of the pod used by the task.


% ─────────────────────────────────────────────────────────────────────────────
\section{Code Structure}

The project is organised on four logical levels: \texttt{core} (shared entities and data structures), \texttt{sim} (DES engine and event handlers), \texttt{opt} (assignment policies) and \texttt{stat} (KPIs cellection management).

\begin{table}[htbp]
\begin{tabularx}{\textwidth}{l X}
\toprule
\textbf{Path} & \textbf{Role} \\
\midrule
\texttt{run\_simulation.py}          & Entry point: configures logging, builds the warehouse, starts the simulation \\
\texttt{experiments.csv}             & Centralised numerical parameters \\
\texttt{scripts/core/enums.py}       & Enumerations for discrete states and event types \\
\texttt{scripts/core/entities.py}    & Dataclasses for \texttt{Order}, \texttt{Pod}, \texttt{Robot}, \texttt{Workstation} and \texttt{Task}. \\
\texttt{scripts/core/queues.py}      & Generic priority queue (min-heap with O(1) lookup) \\
\texttt{scripts/core/warehouse.py}   & Physical layout and generation, distance utilities \\
\texttt{scripts/sim/Simulator.py}    & Main DES engine: clock, event loop, dispatch table \\
\texttt{scripts/sim/event\_handler.py} & Handler function for each event type \\
\texttt{scripts/stat/StatManager.py} & Main collector of KPIs \\
\texttt{scripts/stat/core.py}        & Dataclasses to support the collection of KPIs \\
\texttt{scripts/opt/policies.py}     & Greedy assignment policies \\
\texttt{scripts/opt/OptManager.py}   & The class \texttt{OptManager} controls the optimisation process: precomputes the time-space network and converts binary variables in \texttt{Task} objects\\
\texttt{scripts/opt/local\_search\_stage1.py} & Definition of the local search algorithm used by the otimization solver for the stage  1\\
\texttt{scripts/opt/local\_search\_stage2.py} & Definition of the local search algorithm used by the otimization solver for the stage  2\\
\texttt{scripts/opt/stage2\_data.py}  & Dataclass for the second stage of the optimization process \\
\texttt{scripts/opt/utils.py}  & Here is defined a function that converts the binary variables into \texttt{Task} used by the simulator. \\
\bottomrule
\end{tabularx}
\caption{Folder and file structure of the project.}
\end{table}


% ─────────────────────────────────────────────────────────────────────────────
\newpage
\section{Input Parameters}

The following table lists the simulator configuration parameters.
The parameters that define order generation and pod processing times at workstations have a domain defined by \cite{Barnhart2024}\footnote{GitHub repository \url{https://github.com/alexschmid3/warehouse-task-assignment}.}.

\begin{table}[htbp]
\begin{longtable}{p{5cm} p{7.5cm}}
\toprule
\textbf{Parameter} & \textbf{Description} \\
\midrule
\endfirsthead
\toprule
\textbf{Parameter} & \textbf{Description} \\
\midrule
\endhead
\multicolumn{2}{l}{\textit{Layout and capacity}} \\[2pt]
\texttt{NUM\_PODS}                   & Total number of pods (\texttt{GRID\_ROWS} $\times$ \texttt{GRID\_COLS}) \\
\texttt{NUM\_SKUS}                   & Number of distinct SKUs (fixed in such way that in average every pod has 20 skus) \\
\texttt{NUM\_SKUS\_PER\_POD}         & Target number of distinct SKUs for each pod \\
\texttt{NUM\_ROBOTS}                 & Number of autonomous mobile robots \\
\texttt{NUM\_WORKSTATIONS}           & Number of picking stations \\
\texttt{GRID\_ROWS / GRID\_COLS}     & Pod grid dimensions \\
\texttt{WS\_ORDER\_CAPACITY}         & Maximum number of open orders per station \\
\texttt{WS\_RELEASED\_TASK\_CAPACITY}& Number of tasks released per workstation \\[6pt]
\multicolumn{2}{l}{\textit{Operational times}} \\[2pt]
\texttt{ROBOT\_SPEED}                & Robot speed \\
\texttt{POD\_PROCESS\_TIME}          & Processing time of a pod at the workstation \\
\texttt{ITEM\_PROCESS\_TIME}         & Picking time of an item at the workstation \\[6pt]
\multicolumn{2}{l}{\textit{Order generation}} \\[2pt]
\texttt{INTERARRIVAL\_TIME\_ORDER}   & Mean order inter-arrival time \\
\texttt{PROB\_1\_ITEM\_ORDER}        & Probability that an order contains a single item \\
\texttt{GEO\_DIST\_PARAM\_ORDER}     & Geometric parameter for multi-item orders \\[6pt]
\multicolumn{2}{l}{\textit{Simulation control}} \\[2pt]
\texttt{TIME\_HORIZON}              & Total simulation duration  \\
\texttt{DELTA\_T\_OPT}              & Optimiser call interval \\
\bottomrule
\end{longtable}
\caption{Numerical parameters of the simulator divided by functionality.}
\end{table}

The order generation algorithm is reported here for completeness, as it operationalizes the stochastic rules defined by the parameters table.

\subsection{Order Generation}
Order generation follows the model of Barnhart et al \cite{Barnhart2024}:
\begin{algorithm}
\caption{Order Generation}
\begin{algorithmic}[1]
\State Sample $u \sim \text{Uniform}(0,1)$
\If{$u < \texttt{PROB\_1\_ITEM\_ORDER}$}
    \State $order\_size \gets 1$
\Else
    \State Sample $g \sim \text{Geometric}(p)$ with $p = \texttt{GEO\_DIST\_PARAM\_ORDER}$ until $g \le 15$
    \State $order\_size \gets g + 1$
\EndIf

\For{$i = 0$ to $order\_size - 1$}
    \State Sample SKU$_i$ from a truncated normal distribution centered in $N/2$ and with standard deviation $N/6$ \Comment{$N$ is the total number of unique SKUs in the warehouse}
\EndFor
\end{algorithmic}
\end{algorithm}

% ─────────────────────────────────────────────────────────────────────────────
\newpage \clearpage
\section{Domain Entities}

Entities are implemented as Python dataclasses and represent the state of the system during the simulation. The main ones are: \texttt{Order}, \texttt{Pod}, \texttt{Robot}, \texttt{Workstation}, \texttt{Task} and \texttt{Visit}.

\subsection{Order}

Tracks a customer order from arrival to completion. 

\vspace{0.4em}
\begin{tabularx}{\textwidth}{l l X}
\toprule
\textbf{Field} & \textbf{Type} & \textbf{Description} \\
\midrule
\texttt{order\_id}       & \texttt{int}         & Unique identifier \\
\texttt{arrival\_time}   & \texttt{float}       & Order arrival time \\
\texttt{order\_size}     & \texttt{int}         & Number of items required \\
\texttt{items\_required} & \texttt{list[int]}   & Full list of required items \\
\texttt{items\_pending}  & \texttt{list[int]}   & Items not yet picked (decremented as items are picked)\\
\texttt{workstation\_id} & \texttt{int | None}  & Assigned workstation (\texttt{None} if not yet assigned) \\
\texttt{status}          & \texttt{OrderStatus} & Lifecycle state: \texttt{BACKLOG → WAITING → OPEN → CLOSED} \\
\bottomrule
\end{tabularx}

\subsection{Pod}

Represents a storage pod. When idle, the pod is at its \texttt{storage\_location}; when on a mission, its status becomes \texttt{BUSY}, and the location is not known precisely as routing is not considered.

\vspace{0.4em}
\begin{tabularx}{\textwidth}{l l X}
\toprule
\textbf{Field} & \textbf{Type} & \textbf{Description} \\
\midrule
\texttt{pod\_id}           & \texttt{int}            & Unique identifier \\
\texttt{storage\_location} & \texttt{int}            & Identinfier of the storage cell \\
\texttt{items}             & \texttt{list[int]}      & SKUs contained in the pod \\
\texttt{status}            & \texttt{PodStatus}      & \texttt{IDLE} or \texttt{BUSY} \\
\bottomrule
\end{tabularx}

\subsection{Robot}

Each robot is always in one of two states: \texttt{IDLE} (available) or \texttt{BUSY} (executing a task). 

\vspace{0.4em}
\begin{tabularx}{\textwidth}{l l X}
\toprule
\textbf{Field} & \textbf{Type} & \textbf{Description} \\
\midrule
\texttt{robot\_id}         & \texttt{int}            & Unique identifier \\
\texttt{position}          & \texttt{id}             & Current position in the grid in terms of cell id (the information is known only when the robot is not moving)\\
\texttt{status}            & \texttt{RobotStatus}    & \texttt{IDLE} or \texttt{BUSY} \\
\bottomrule
\end{tabularx}

\subsection{Workstation}

The workstation manages two separate buffers: one for pods (\texttt{picking\_buffer}) and one for orders waiting to be opened (\texttt{order\_buffer}). Note that the \texttt{status} refers to the picking operation, not to whether any orders are currently open at the workstation.

\vspace{0.4em}
\begin{tabularx}{\textwidth}{l l X}
\toprule
\textbf{Field} & \textbf{Type} & \textbf{Description} \\
\midrule
\texttt{workstation\_id}         & \texttt{int}                            & Unique identifier \\
\texttt{order\_capacity}         & \texttt{int}                            & Maximum simultaneously open orders \\
\texttt{released\_task\_capacity}& \texttt{int}                      & Maximum released tasks per workstation \\
\texttt{position}                & \texttt{int}                      & Cell identifier of the position \\
\texttt{pod\_process\_time}      & \texttt{float}                          & Pod processing time at a workstation \\
\texttt{item\_process\_time}     & \texttt{float}                          & Time for a picking operation at a workstation \\
\texttt{opened\_orders}          & \texttt{set[int]}                       & Orders currently being processed \\
\texttt{order\_buffer}           & \texttt{list[int]}                      & Orders waiting to be opened \\
\texttt{picking\_buffer}         & \texttt{list[int]}                      & Pods waiting (excluding the active one) for a picking operation\\
\texttt{active\_tasks}           & \texttt{set[int]}                       & IDs of released tasks allocated to a robot \\
\texttt{released\_tasks}         & \texttt{set[int]}                       & IDs of released tasks not yet allocated \\
\texttt{status}                  & \texttt{WorkstationPickingStatus}       & \texttt{IDLE} or \texttt{BUSY} \\
\bottomrule
\end{tabularx}

\subsection{Task and Visit}

A \texttt{Task} is a mission assigned to a pod-robot pair, composed of an ordered sequence of stops (\texttt{stops}). Each stop is a \texttt{Visit}, which specifies the destination workstation, the orders served (\texttt{orders}), and the items to be picked (\texttt{items}). The \texttt{priority} field of the task determines the dispatching order in the queue.

\subsection{Enumerations}

\texttt{OrderStatus}, \texttt{RobotStatus}, \texttt{PodStatus} and \texttt{WorkstationPickingStatus} define the discrete states of the entities, making state transitions explicit in the simulation logic.

% ─────────────────────────────────────────────────────────────────────────────
\subsection{Warehouse}

The \texttt{Warehouse} class collects all information and state variables of the system entities. In particular, its attributes include the lists of pods, workstations, and robots.

This class also provides methods to generate the warehouse layout based on the number of pods along the $x$-axis, the number of pods along the $y$-axis, the number of workstations and the number of robots.

\subsubsection{Warehouse Layout}

The warehouse layout is generated from \texttt{grid\_rows} and \texttt{grid\_cols}. Pods are arranged in a rectangular grid; after every two columns of pods, an aisle of width two cells is inserted to allow robot movement. The resulting physical dimensions are:

\[
X = \text{grid\_rows} + 2\left\lfloor\frac{\text{grid\_rows}-1}{2}\right\rfloor + 2 \cdot \text{MARGIN} - 1,
\qquad
Y = \text{grid\_cols} + 2 \cdot \text{MARGIN} - 1
\]

Workstations are placed along the perimeter, preferably on the bottom side in a symmetric configuration. Robots are initialized at random non-overlapping positions within the warehouse.  
The travel time between two cells is computed as the sum of two components: a nominal travel time, obtained by dividing the Manhattan distance by \texttt{robot\_speed}, and a positive stochastic component accounting for potential delays due to congestion or other effects.

\begin{figure}[htbp]
\centering
\includegraphics[width=\textwidth]{images/warehouse_layout.png}
\caption{Example layout created with \texttt{GRID\_ROWS} = 10, \texttt{GRID\_COLS} = 10, \texttt{NUM\_ROBOTS} = 15 and \texttt{NUM\_WORKSTATION} = 9. \\ Black cells are the \texttt{storage\_location} of each pod, red ones belong to workstations, and blue ones are the initial positions of the robots.}
\end{figure}

% ─────────────────────────────────────────────────────────────────────────────
\newpage \clearpage
\section{Simulator Engine}

The simulation engine is split across three Python objects to cleanly separate immutable configuration from mutable per-run state.

% ── SimulatorConfig ──────────────────────────────────────────────────────────
\subsection{\texttt{SimulatorConfig} - immutable parameters}

A dataclass that holds every parameter that does not change across replicas. It is constructed once and shared by all \texttt{run()} calls on the same \texttt{Simulator} instance.

\vspace{0.4em}
\begin{tabularx}{\textwidth}{l l X}
\toprule
\textbf{Attribute} & \textbf{Type} & \textbf{Description} \\
\midrule
\texttt{order\_gen\_config}    & \texttt{list[float]}  & \texttt{[interarrival\_time, prob\_1\_item, geo\_param]}, parameters for order arrival generation \\
\texttt{time\_horizone}        & \texttt{float}        & Time horizon of the simulator in seconds.. \\
\texttt{warm\_up}              & \texttt{float}        & Seconds discarded from statistics at run start. Default \texttt{0}. \\
\texttt{optimization\_enabled} & \texttt{bool}         & If \texttt{False}, greedy policies are adopted; if \texttt{True}, optimiser is adopted \\
\texttt{optimization\_interval} & \texttt{float}        & Time interval between two optimization in seconds \\
\bottomrule
\end{tabularx}

% ── SimulatorState ───────────────────────────────────────────────────────────
\subsection{\texttt{SimulatorState} - mutable per-run state}

A plain dataclass holding all quantities that change during a simulation run. A fresh \texttt{SimulatorState} is constructed at the beginning of every \texttt{run()} call, guaranteeing that consecutive calls on the same \texttt{Simulator} object produce fully independent replicas with no shared mutable state.

\vspace{0.4em}
\begin{tabularx}{\textwidth}{l l X}
\toprule
\textbf{Attribute} & \textbf{Type} & \textbf{Description} \\
\midrule
\texttt{current\_time}      & \texttt{float}                          & Current simulation time (seconds) \\
\texttt{warehouse}          & \texttt{Warehouse}                      & Physical environment and live state of all entities \\
\texttt{future\_events}     & \texttt{PriorityQueue[Event]}           & Future events ordered by timestamp; \texttt{RELEASE\_TASK} breaks ties first \\
\texttt{orders\_in\_system} & \texttt{PriorityQueue[Order]}           & Received orders ordered by arrival time; \texttt{BACKLOG} orders have priority \\
\texttt{released\_tasks}    & \texttt{PriorityQueue[Task]}            & Tasks ready for robot assignment \\
\texttt{active\_tasks}      & \texttt{dict[int, Task]}                & Executing tasks, keyed by \texttt{task\_id} \\
\texttt{orders\_counter}    & \texttt{int}                            & Monotonic counter for order IDs \\
\texttt{task\_counter}      & \texttt{int}                            & Monotonic counter for task IDs \\
\texttt{event\_count}       & \texttt{int}                            & Total events processed (for logging) \\
\texttt{closed\_orders}     & \texttt{int}                            & Completed orders (for logging) \\
\bottomrule
\end{tabularx}

% ── Simulator ────────────────────────────────────────────────────────────────
\subsection{\texttt{Simulator}}

The \texttt{Simulator} class owns the immutable configuration, the random generator and the \texttt{StatManager}. Its constructor accepts a \texttt{warehouse\_factory}, a zero-argument callable that returns a fresh \texttt{Warehouse}, instead of a \texttt{Warehouse} instance directly. This design choice is intentional: calling \texttt{run()} multiple times in the main script produces independent replicas automatically, because each call rebuilds the full \texttt{SimulatorState} (including a new \texttt{Warehouse}) from scratch before processing any event.

\vspace{0.4em}
\begin{tabularx}{\textwidth}{l l X}
\toprule
\textbf{Attribute} & \textbf{Type} & \textbf{Description} \\
\midrule
\texttt{RANDOM\_GENERATOR}    & \texttt{numpy.Generator}          & Random generator for reproducibility \\
\texttt{config}               & \texttt{SimulatorConfig}          & Immutable run parameters \\
\texttt{state}                & \texttt{SimulatorState}           & Mutable state of the current run (rebuilt on every \texttt{run()} call) \\
\texttt{STAT\_MANAGER}        & \texttt{StatManager}              & KPI collector (also rebuilt on every \texttt{run()} call) \\
\texttt{\_warehouse\_factory} & \texttt{Callable[[], Warehouse]}  & Factory invoked once per \texttt{run()} to produce a pristine warehouse \\
\texttt{OPT\_MANAGER}         & \texttt{OptManager}               & Optimization manager (also rebuilt on every \texttt{run()} call), which precomputes the time-space network and provides the interface to the optimizer. \\
\bottomrule
\end{tabularx}

\vspace{0.4em}
The main event loop inside \texttt{run()} follows the classical DES structure:
\begin{enumerate}[noitemsep]
    \item Rebuild \texttt{state} and \texttt{STAT\_MANAGER} from a fresh warehouse.
    \item Push the first \texttt{ARRIVAL\_ORDER} event at $t \approx 0$.
    \item If optimization is enabled, push the first \texttt{RUN\_OPTIMIZER} event at  $t \approx 1 min$.
    \item Pop the event with the smallest timestamp from \texttt{state.future\_events}.
    \item Advance \texttt{state.current\_time} to the event time.
    \item Dispatch to the corresponding handler via the dispatch table.
    \item Repeat until the queue is empty or \texttt{TIME\_HORIZON} is reached.
    \item Write and reset statistics via \texttt{STAT\_MANAGER}.
\end{enumerate}

To be precise, the first \texttt{ARRIVAL\_ORDER} event will make arrive $100$ orders simultaneously to mimic a initial situation of a already working warehouse.


% ─────────────────────────────────────────────────────────────────────────────
\subsection{Event Handlers}

Each \texttt{EventType} is associated with a handler function with signature \texttt{handler(event, state, sim)}, where \texttt{state} is the \texttt{SimulatorState} holding all mutable queues and counters, and \texttt{sim} provides access to the immutable configuration (\texttt{sim.config}) and the random generator (\texttt{sim.RANDOM\_GENERATOR}). The table summarises the current implementation status.

\begin{table}[htbp]
\begin{tabularx}{\textwidth}{p{3cm} p{4cm} X}
\toprule
\textbf{Event} & \textbf{\texttt{Event.info}} & \textbf{Responsibility} \\
\midrule
\texttt{ARRIVAL\_ORDER}    
& \texttt{None} 
& Generates a new order, store it in \texttt{state.orders\_in\_system} and schedules the next order arrival; without the optimiser, assigns it directly to a workstation. \\

\texttt{OPEN\_ORDER}       
& \texttt{Order} 
& Sets the order to \texttt{OPEN}, inserts it into \texttt{opened\_orders} and, if possible, generates the associated tasks and triggers \texttt{RELEASE\_TASK}. \\

\texttt{RELEASE\_TASK}     
& \texttt{Task}
& Releases tasks, if optimizer is enabled all \texttt{RELEASE\_TASk} events will be scheduled at the end of the optimization; otherwise events will be scheduled when tasks are designed. \\

\texttt{START\_TASK}       
& \texttt{None} 
& Verify that among the released tasks there is at least one that can be started (i.e., whose corresponding pod is idle). If so, assign an idle robot (if available) to it and set both the robot and the pod to \texttt{BUSY}. \\

\texttt{ARRIVAL\_POD\_WST} 
& \texttt{Task} 
& Handles pod arrival at the workstation: inserts it into the \texttt{picking\_buffer} and triggers \texttt{START\_PICKING} if the workstation is idle. \\

\texttt{START\_PICKING}    
& \texttt{Task} 
& Starts picking, sets the workstation to \texttt{BUSY} and schedules \texttt{END\_PICKING} according to the picking process times. \\

\texttt{END\_PICKING}      
& \texttt{Task} 
& Completes the picking operation by updating the \texttt{items\_pending} field of the involved orders. If necessary, it triggers \texttt{CLOSE\_ORDER}, and schedules new pod movements or their return to storage based on the remaining stops in the task. Since the number of released tasks for this workstation has been decremented, new tasks can now be created and released.\\

\texttt{CLOSE\_ORDER}      
& \texttt{Order} 
& Closes the order when completed, frees a slot at the workstation and trigger \texttt{OPEN\_ORDER} if possible. \\

\texttt{RETURN\_POD}       
& \texttt{Task} 
& Returns the pod to its \texttt{storage\_location}, sets robot and pod to \texttt{IDLE} and enables new tasks to start. \\

\texttt{RUN\_OPTIMIZER}    
& \texttt{None} 
& Invokes the optimiser to assign backlog orders and generate tasks. To mantain coherence between simulation and optimization, active tasks are considered as completed while released tasks not yet allocated are overwritten. \\
\bottomrule
\end{tabularx}
\end{table}

Note that once a task is extracted from the priority queue in \texttt{START\_TASK}, it is then passed along between events as input and moves through the simulation.

In the \texttt{Simulator} class, when the \texttt{future\_events} queue is initialized, ties between simultaneous events are resolved by prioritizing \texttt{RELEASE\_TASK} events. This is because the greedy policy used for task generation must consider the full set of released tasks in order to avoid producing redundant tasks, which could lead to inconsistencies in the simulation. For the same reason, at the end of the \texttt{END\_PICKING} event, if a new order is opened at the workstation, no new tasks are designed or released, as this will be handled during the \texttt{OPEN\_ORDER} process.  
The reason behind these design choices will become clearer when the greedy policies are introduced.

The following flowcharts illustrate the dependencies between events.

\begin{figure}[htbp]
\centering\includegraphics[width=\textwidth]{images/Diagramma_ERG_Policy.pdf}
\caption{Graph showing the relationship between events when \texttt{optimization\_enabled = False}. Arrows indicate which event will be scheduled next.}
\end{figure}

\begin{figure}[htbp]
\centering\includegraphics[width=\textwidth]{images/Diagramma_ERG_Opt.pdf}
\caption{Graph showing the relationship between events when \texttt{optimization\_enabled = True}. Arrows indicate which event will be scheduled next.}
\end{figure}


% ─────────────────────────────────────────────────────────────────────────────
\newpage \clearpage
\section{KPIs Collection}
The main objective of this simulator is to evaluate two different optimization approaches in a controlled and reproducible environment; hence, relevant statistics must be collected.

The collected statistics are grouped into three categories:
\begin{itemize}
    \item \texttt{ResourceTracker}: this type of statistic measures the state frequency of a given entity over time. It is used, for example, for both robots and workstations, and captures respectively the robot utilization rate and the workstation picking time ratio.
    \item \texttt{OrderFlowTracker}: this statistic captures the average order flow time, grouped by order size, from the \texttt{BACKLOG} state to the \texttt{CLOSED} state.
    \item \texttt{TimeWeightedMeanTracker}: this type of statistic is used to compute time-averaged KPIs, such as the average number of open orders per workstation or the average number of pods simultaneously out of their storage location.
    \item \texttt{Counter}: this type of statistic is used to count occurencies of something; in particular, it is used to track the throughtput of the system.
\end{itemize}




\newpage
\section{Optimization Approaches}
\subsection{Greedy Policies}

\subsubsection{Assignment Order-Workstation}
The assignment order-workstation is performed whenever an order arrives in the system.  
The greedy policy implemented is the following\footnote{Vorrei migliorarla cercando di creare dei cluster in base alle sku required dagli order alle workstation}:
\begin{algorithm}
\caption{Assignment order-workstation}
\begin{algorithmic}[1]
\State Select \texttt{workstation\_id} as the workstation minimizing the sum of queued and currently open orders
\end{algorithmic}
\end{algorithm}

In this policy, the sequencing of orders at a given workstation is neglected, as their priority is solely given by their arrival time.

\subsubsection{Task Design and Sequencing}
Task design is performed greedily, considering only the orders currently open at a given workstation, with the objective of maximizing pile-on that is, minimizing the number of pods required to complete all open orders at the workstation.  
During the algorithm, all released tasks associated with the given workstation that have not yet been dispatched are redesigned to account for newly required items. This is why the realising of already designed tasks must be prioritized in order to avoid inconsistencies.

In practice, task sequencing is neglected, as priority is based exclusively on the time at which tasks are designed and released.

\begin{algorithm}
\caption{Task Design}
\begin{algorithmic}[1]
\State Compute \texttt{uncovered\_skus} as the set of SKUs required by open orders at the workstation that are not yet covered by active tasks
\If{\texttt{uncovered\_skus} is empty}
    \State \Return
\EndIf

\While{the number of designed tasks is below the maximum number of tasks that can be generated \textbf{and} \texttt{uncovered\_skus} is not empty}
    \State Rank pods according to the number of items in \texttt{uncovered\_skus} they can serve (breaking ties by distance)
    \State Select the highest-ranked pod
    \State Define a \texttt{Visit} based on the subset of orders that benefit from the selected pod
    \State Create a task containing the defined \texttt{Visit}
\EndWhile
\State \Return Tasks designed
\end{algorithmic}
\end{algorithm}

Some observations have to be made. 
First, this policy does not allow tasks with multiple visits to be defined, as it does not have global information about all workstations, but only about the one that triggers the policy.
Second, in this policy the while loop continues until the number of designed tasks reaches the maximum number of tasks that can be generated. This is due to the fact that the number of released and active tasks per workstation is limited by \texttt{WS\_RELEASED\_TASK\_CAPACITY}. This design choice is intended to prevent overcrowding near the workstation in scenarios where the number of robots significantly exceeds the system's processing capacity (a large fleet of robots serving a small number of slow workstations, for istance), which would otherwise lead to unnecessary congestion.


\subsection{Global optimization model}

The model is a static optimization problem split into two stages: the first stage assigns orders to workstations and pods to orders, while the second stage designs tasks and schedules them over a time-space network. The two stages are connected by the decision variables of the first stage, which are used as input parameters for the second stage.

\subsubsection{Time-space network}
The time-space network is constructed at the beginning of each simulation run, when the \texttt{OptManager} class is instantiated. The network is built by considering arcs of the form
\[
\left((s_s, t_s), (s_d, t_d)\right), \qquad t_s < t_d,
\]
where each time-space tuple $(s,t)$ represents a node characterized by a spatial position and a temporal coordinate. The spatial position corresponds to the cell identifier within the warehouse grid, while the time coordinate is defined relative to the considered optimization run.

Not all possible arcs are included in the network, as this would significantly increase its size. Specifically, only two classes of arcs are considered:
\begin{itemize}
    \item \emph{travelling arcs}, namely arcs for which either the source location or the destination location corresponds to a workstation, and such that $t_d - t_s$ is sufficient for the pod to reach the workstation under deterministic travel times; 
    \item \emph{idle arcs}, namely arcs satisfying
    \[
    s_s = s_d, \qquad t_d = t_s + 1,
    \]
    where $s_s$ corresponds to a pod storage location.
\end{itemize}

\subsubsection{First Stage: Order-Workstation and Order-Item-Pod Assignments}

The first stage assigns orders to workstations and items to pods. The decision
variables are:
\begin{itemize}[noitemsep]
    \item $z_{mw} = 1$ if order $m$ is assigned to workstation $w$, $0$ otherwise.
    \item $x_{imp} = 1$ if item $i$ of order $m$ is retrieved from pod $p$, $0$ otherwise.
    \item $y_{wp} = 1$ if pod $p$ is assigned to workstation $w$, $0$ otherwise.
\end{itemize}

In the implementation, $x_{imp}$ is stored with two indices only, since not all items $i$ are required by order $m$; a bijective map $(i,m)\to im$ tracks the correspondence between the two representations.

\begin{align}
    \min \quad
    & \sum_{w \in \mathcal{W}}\sum_{p \in \mathcal{P}} y_{wp} \\[4pt]
    \text{s.t.} \quad
    & \sum_{w \in \mathcal{W}} z_{mw} = 1
      \qquad \forall m \in \mathcal{M} \label{M1_Constr1} \\
    & \sum_{p \in \mathcal{P}_i} x_{imp} = 1
      \qquad \forall m \in \mathcal{M},\, i \in \mathcal{I}_m \label{M1_Constr2} \\
    & y_{wp} \ge x_{imp} + z_{mw} - 1
      \qquad \forall w \in \mathcal{W},\, m \in \mathcal{M},\,
              i \in \mathcal{I}_m,\, p \in \mathcal{P}_i \label{M1_Constr3} \\
    & \sum_{m \in \mathcal{M}} |\mathcal{I}_m|\, z_{mw} \le \overline{I}
      \qquad \forall w \in \mathcal{W} \label{M1_Constr4} \\
    & \sum_{m \in \mathcal{M}} |\mathcal{I}_m|\, z_{mw} \ge \underline{I}
      \qquad \forall w \in \mathcal{W} \label{M1_Constr5} \\
    & z_{mw} = 1
      \qquad \forall m \in \mathcal{M}_0 \label{M1_InitialCond}
\end{align}

where $\mathcal{M}_0 \subset \mathcal{M}$ is the subset of orders whose workstation is already fixed by the current warehouse state. The objective minimises the total number of (workstation, pod) pairs, which is a proxy for the number of pod movements. Constraints \eqref{M1_Constr4}--\eqref{M1_Constr5} enforce workload balance; the bounds are set to $\overline{I} = 1.15\,\frac{\sum_m |\mathcal{I}_m|}{|\mathcal{W}|}$ and $\underline{I} = 0.85\,\frac{\sum_m |\mathcal{I}_m|}{|\mathcal{W}|}$.

The problem is solved via the local search described in Section~\ref{sss:ls_stage1}; its output, the assignments $z$ and $x$, is passed as input to the second stage.


\subsubsection{Second Stage: Task Design and Scheduling}

Given the assignments from Stage 1, the second stage schedules pod movements over the time-space network. The decision variables are:
\begin{itemize}[noitemsep]
    \item $x_{imt} = 1$ if item $i$ of order $m$ has been picked by time $t$,
          $0$ otherwise.
    \item $f_{mt} = 1$ if at least one item of order $m$ has been picked by time
          $t-1$, $0$ otherwise.
    \item $g_{mt} = 1$ if order $m$ is completed by time $t-1$, $0$ otherwise.
    \item $v_{mt} = 1$ if order $m$ is open at time $t$, $0$ otherwise.
    \item $y_{pa} = 1$ if arc $a$ in the time-space network is selected for pod
          $p$, $0$ otherwise.
\end{itemize}

As in Stage 1, $x_{imt}$ is stored with two indices via the same bijective map.
The variable $y_{pa}$ is defined only for pods assigned to at least one workstation in Stage 1.

\begin{align}
    \max \quad
    & \sum_{m \in \mathcal{M}} x_{im,T}
      - \sum_{m \in \mathcal{M}} \sum_{t \in \mathcal{T}} \beta_m(t)\,(1 - g_{mt})
      \label{M2_Obj} \\[4pt]
    \text{s.t.} \quad
    & \sum_{m \in \mathcal{M}(w)} v_{mt} \le C_w
      \qquad \forall w \in \mathcal{W},\, t \in \mathcal{T}
      \label{M2_Constr1} \\
    & \sum_{m \in \mathcal{M}(w)} \sum_{i \in \mathcal{I}_m}
      \delta^{\text{item}} (x_{imt} - x_{im,t-1})
      + \sum_{p \in \mathcal{P}} \sum_{a \in \mathcal{A}^+_{N(w,t)}}
      \delta^{\text{pod}}\, y_{pa}
      \le \Delta
      \qquad \forall w \in \mathcal{W},\, t \in \mathcal{T} \setminus \{0\}
      \label{M2_Constr2} \\
    & \sum_{a \in \mathcal{A}^+_{N_p^0}} y_{pa} = 1
      \qquad \forall p \in \mathcal{P}
      \label{M2_Constr3} \\
    & \sum_{a \in \mathcal{A}^-_{n}} y_{pa}
      - \sum_{a \in \mathcal{A}^+_{n}} y_{pa} = 0
      \qquad \forall p \in \mathcal{P},\;
              n \in \mathcal{N} \setminus \left(\mathcal{N}_{\text{end}}
              \cup \{N_p^0\}\right)
      \label{M2_Constr4} \\
    & x_{imt} - x_{im,t-1}
      \le \sum_{a \in \mathcal{A}^+_{N(w,t)}} y_{pa}
      \qquad \forall w,\, m \in \mathcal{M}(w),\, i \in \mathcal{I}_m,\,
              t \in \mathcal{T} \setminus \{0\} \\
    & x_{imt} \ge x_{im,t-1}
      \qquad \forall m,\, i \in \mathcal{I}_m,\, t \in \mathcal{T} \setminus \{0\} \\
    & v_{mt} = f_{mt} - g_{mt}
      \qquad \forall m,\, t \in \mathcal{T} \setminus \{0\} \\
    & f_{mt} \ge x_{imt},\quad g_{mt} \le x_{im,t-1}
      \qquad \forall m,\, i \in \mathcal{I}_m,\, t \in \mathcal{T} \\
    & g_{m,t+1} \ge \sum_{i \in \mathcal{I}_m} x_{imt} - (|\mathcal{I}_m| - 1)
      \qquad \forall m,\, t \in \mathcal{T} \\
    & v_{mt} \ge v_{m,t-1} - g_{mt}
      \qquad \forall m,\, t \in \mathcal{T}
      \label{Enforcing_IC} \\
    & v_{m,0} = 0
      \qquad \forall m \in \mathcal{M}_0
      \label{IC} \\
    & f_{mt} \ge f_{m,t-1},\quad g_{mt} \ge g_{m,t-1}
      \qquad \forall m,\, t \in \mathcal{T} \\
    & \sum_{p \in \mathcal{P}} \sum_{a \in \mathcal{A}^{\text{mov}}_p(t)} y_{pa}
      \le |\mathcal{R}|
      \qquad \forall t \in \mathcal{T}
      \label{max_active_pods}
\end{align}

The objective \eqref{M2_Obj} maximises throughput and penalises uncompleted orders with a weight $\beta_m(t)$ proportional to the time elapsed since order $m$ arrived, preventing orders from being left indefinitely in the backlog.
Constraints \eqref{IC}--\eqref{Enforcing_IC} enforce initial conditions for orders already open at workstations, for which $f_{mt}$ is set to $1$ regardless of whether any item has been picked. Constraint \eqref{max_active_pods} bounds the
number of pods simultaneously out of storage, where
\[
\mathcal{A}^{\text{mov}}_p(t) =
\left\{a \in \mathcal{A}_p \;\middle|\;
\ell^+(a) \neq s_p,\quad t^-(a) < t \le t^+(a)\right\}
\]
and $t^-(a)$, $t^+(a)$ denote the departure and arrival times of arc $a$.
The problem is solved via the local search described in Section~\ref{sss:ls_stage2}.


\subsubsection{Local Search Stage 1}
\label{sss:ls_stage1}

The Stage-1 local search minimises $\sum_{w,p} y_{wp}$. At each iteration, a random subset of at most $600$ candidate moves is generated from three types:

\begin{itemize}[noitemsep]
    \item \textbf{Swap.} The workstation assignments of two free orders are exchanged.
    \item \textbf{Move-to.} A single free order is reassigned to a different workstation.
    \item \textbf{Re-pod.} A single item is reassigned to a different pod stocking its SKU.
\end{itemize}

After each move, $y$ is recomputed deterministically from $x$ and $z$. Infeasible candidates are discarded; among the remaining ones, the move yielding the lowest objective is accepted. Ties are accepted to allow plateau traversal. Orders already open at a workstation are never moved. 

\subsubsection{Local Search Stage 2}
\label{sss:ls_stage2}

\paragraph{Initial solution.}
The local search requires a feasible starting point for $x_{imt}$, built by a dedicated construction heuristic. For each workstation $w$, orders are partitioned into batches of at most $C_w$ orders. Batches are grown from a seed order by iteratively adding the candidate that maximises pod sharing and minimises new pods introduced; underfull batches are subsequently merged via best-fit on the same score. 
Pod arrivals are then scheduled tentatively, accounting for pod chaining (the departure location is either the pod's storage cell or the last workstation it was committed to). 
A batch is committed only if all its required pods found a feasible slot; uncommitted orders are retried individually.

The construction heuristic is summarised in Algorithm~\ref{alg:initial_solution}.

\begin{algorithm}
\caption{Initial Solution Construction}
\label{alg:initial_solution}
\begin{algorithmic}[1]
\For{each workstation $w$}
    \State Partition orders into batches via greedy pod-sharing score
    \State Merge underfull batches with best-fit
    \State Prepend opened orders as first batch
    \For{each batch $b$}
        \State Sort pods by order count (desc) and earliest pick time
        \For{each pod $p$ in sorted order}
            \State $t^* \gets$ \Call{FindSlot}{$p$, earliest feasible time, travel time}
            \If{$t^* \neq \texttt{None}$} register tentative pick at $t^*$ \EndIf
        \EndFor
        \If{all pods of at least one order found a slot}
            \State Commit picks; set $x_{im,\,t^*:} = 1$ for each item $i$, order $m$
        \Else
            \State Add orders of the batch to the retry pool
        \EndIf
    \EndFor
    \State Retry uncommitted orders individually
\EndFor
\State \Return $x$
\end{algorithmic}
\end{algorithm}

\paragraph{Neighbourhood and move types.}
The search perturbs only $x_{imt}$; all other variables are derived from $x$. Up to $150$ candidate moves are sampled per iteration, allocated $50\%$--$30\%$--$20\%$ across the three classes:

\begin{itemize}[noitemsep]
    \item \textbf{Item shift:} The pick time of one item, or of a small random subset of items, is shifted earlier by a fixed number of time steps.
    \item \textbf{Order shift:} All items of an order are shifted earlier by a fixed number of time steps.
    \item \textbf{Order swap:} Two orders at the same workstation exchange their position in the picking sequence.
\end{itemize}

\paragraph{Two-phase feasibility check.}
To avoid rebuilding $y$ at every candidate evaluation, feasibility is verified in two stages. First, a cheap $x$-only check verifies the constraints that do not involve pod routing (monotonicity of $x$, capacity \eqref{M2_Constr1}, consistency of $f$, $g$, $v$), which are updated incrementally by recomputing only the rows of affected orders. Only for the best $x$-feasible candidate per iteration are the affected pod rows of $y$ rebuilt and the full constraint check
\eqref{M2_Constr3}--\eqref{max_active_pods} performed. If this check fails, the second-best candidate is tried before discarding the iteration. 

\subsubsection{From Binary Variables to Simulator Objects}

The function contained in \texttt{scripts/opt/utils.py} translates the binary solutions $(x, v, y)$ produced by the two-stage optimiser into the \texttt{Task} and \texttt{Visit} objects seen by the simulator. The conversion proceeds in four steps.

\begin{enumerate}
\item \textbf{Item pick times}
For each relevant $(i, m)$ pair indexed by $im$, the pick time $t^*_{im}$ is the first timestep at which $x_{im,t}$ transitions from $0$ to $1$. Combined with the pod assignment from Stage~1, this defines the lookup $\texttt{pick\_at}[(p, t, w)]$, which collects, for each pod $p$ arriving at workstation $w$ at time $t$, the set of items and orders to be served.

\item \textbf{Pod trajectory reconstruction}
For each pod $p$, the arcs selected by $y_{pa} > 0.5$ are extracted and sorted by departure time. The sequence is walked to identify \emph{trips}: a trip begins when the pod departs from storage and ends when it returns. Each stop corresponding to a workstation visit with a non-empty item set is encoded as a \texttt{Visit}; a trip with at least one valid \texttt{Visit} becomes a \texttt{Task}.

\item \textbf{Priority assignment}
The priority of each task is set to:
\begin{equation}
    \pi = t^*_{\text{first}} - \frac{0.25 \cdot d(p,\, w)}{\Delta_{\text{unit}}}
\end{equation}
where $t^*_{\text{first}}$ is the pick time of the first item at the first stop, $d(p, w)$ is the travel time from the pod's storage cell to the first destination workstation, and $\Delta_{\text{unit}}$ is the optimiser time unit. 

\item \textbf{Sorting and ID assignment}
Tasks are sorted by ascending priority and remapped onto $[0,\, 300]$ via linear rescaling. IDs are assigned sequentially from \texttt{state.task\_counter}, and orders within each workstation are sorted by their start time derived from $v_{mt}$, with ties broken by the priority of the first task serving that order.
\end{enumerate}

Finally, note that the simulator may not respect the optimized task sequence: a task is additionally held back until at least one of its served orders is already open or is first in the order queue at the target workstation. Without this check, a robot could deliver a pod to a workstation with no open order ready to receive it, causing a dead-lock.

\clearpage
\newpage
\printbibliography

\end{document}